% ==========================================================
% PACOTES
% ==========================================================
\usepackage{multicol}
\usepackage{fontspec}
\usepackage{etoolbox}
\usepackage{environ}
\usepackage{float}
\usepackage{witharrows}
\usepackage{eso-pic, pgf-pie}
\usepackage{pgfplots}
\usepackage{cancel}
\usepackage{tikz}
\usepackage{amsmath}
\usepackage{tcolorbox}
\tcbuselibrary{breakable}

\newcommand{\canceldiv}[2]{%
    \cancel{#1}^{\div #2}%
}

\renewcommand{\CancelColor}{\color{red}}

\usepackage{polyglossia}
\setmainlanguage{portuguese}
\setotherlanguage{english}

% Fonte
\setmainfont{TeX Gyre Heros}


% Matemática

\usepackage{unicode-math}

\setmathfont{Latin Modern Math}
% Layout
\usepackage[
    a4paper,
    top=2cm,
    bottom=2cm,
    left=1.5cm,
    right=2.0cm,
    headheight=15pt
]{geometry}

% TikZ
\usepackage{tikz}
\usetikzlibrary{
    calc,
    positioning,
    shapes.geometric,
    shapes.misc,
    arrows.meta,
    backgrounds,
    decorations.pathmorphing,
    fadings
}

% Caixas
\usepackage{tcolorbox}
\tcbuselibrary{
    skins,
    breakable
}

% Listas
\usepackage{enumitem}

% Links
\usepackage[
    colorlinks=true,
    linkcolor=AzulPrincipal,
    urlcolor=Roxo,
    citecolor=AzulPrincipal
]{hyperref}

% Cabeçalho
\usepackage{fancyhdr}

% Controle de títulos
\usepackage{titlesec}

% Espaçamento
\usepackage{setspace}
\setstretch{1.08}

% Imagens
\usepackage{graphicx}

% Tabelas
\usepackage{array}
\usepackage{booktabs}

% ==========================================================
% CONFIGURAÇÕES GERAIS
% ==========================================================

\setlength{\parindent}{0pt}
\setlength{\parskip}{6pt}

\setlist[itemize]{
    leftmargin=*,
    itemsep=2pt
}

\setlist[enumerate]{
    leftmargin=*,
    itemsep=4pt
}

% ==========================================================
% CABEÇALHO E RODAPÉ
% ==========================================================

\pagestyle{fancy}

\fancyhf{}

\fancyhead[LE]{
    \small\color{CinzaTexto}
    Banco de Questões de Matemática
}

\fancyhead[RO]{
    \small\color{CinzaTexto}
    \leftmark
}

\fancyfoot[C]{
    \tikz[baseline=-0.5ex]{
        \node[
            circle,
            fill=AzulPrincipal,
            text=white,
            minimum size=0.65cm,
            inner sep=0pt
        ] {\small\thepage};
    }
}

\renewcommand{\headrulewidth}{0pt}
\renewcommand{\footrulewidth}{0pt}


% ==========================================================
% CAPÍTULOS
% ==========================================================

\titleformat
{\chapter}
[display]
{\normalfont\sffamily}
{}
{0pt}
{%
    \begin{tikzpicture}[overlay]
        \node[
            fill=AzulPrincipal,
            text=white,
            rounded corners=8pt,
            inner xsep=15pt,
            inner ysep=5pt
        ] at (2.4,.3) {\Large\bfseries CAPÍTULO \thechapter};
    \end{tikzpicture}
    \hspace{5cm}
    \Huge\bfseries
}
[
    \vspace{-.9cm}
    \color{Turquesa}
    \rule{\textwidth}{3pt}
]


\titleformat
{name=\chapter,numberless}
[display]
{\normalfont\sffamily}
{}
{0pt}
{
    \Huge\bfseries
}
[
    \vspace{-1.1cm}
    \color{Turquesa}
    \rule{\textwidth}{3pt}
]

\titlespacing*{\chapter}
{0pt}{-20pt}{30pt}

% ==========================================================
% SEÇÕES
% ==========================================================

\titleformat
{\section}
{\Large\bfseries\sffamily\color{AzulPrincipal}}
{\thesection}
{0.7em}
{}

\titleformat
{\subsection}
{\large\bfseries\sffamily\color{Roxo}}
{\thesubsection}
{0.7em}
{}

% =============================================================
% AMBIENTE DE OBSERVAÇÃO OU DICAS
% =============================================================

\newtcolorbox{observacao}[1][Observação]{
    colback=VerdeDestaqueFundo,
    colframe=VerdeDestaque,
    boxrule=0pt,
    leftrule=4pt,
    arc=2pt,
    left=10pt,
    right=10pt,
    top=8pt,
    bottom=8pt,
    enhanced,
    before upper={%
        \textcolor{VerdeDestaque}{\textbf{#1}}
        \par\smallskip
    },
}